% Unit 080 English translation of chapter6.tex lines 1540--1631.
% Source: Methods of Algebra, Volume 2, commit
% 9a5803ff77dd3257484cb177f851a73770a59dd3 (CC BY 4.0).
% stable-unit-id: o014.aljabr2.chapter6.lhs-spectral-sequence-a
% Coverage: Section 6.9, construction and statement of the LHS theorem.

% segment-id: o014.li2.chapter6.lhs-spectral-sequence-a.h001
\section{The Lyndon--Hochschild--Serre spectral sequence}\label{sec:LHS-SS}
\hypertarget{o014.aljabr2.chapter6.lhs-spectral-sequence-a}{ }

% segment-id: o014.li2.chapter6.lhs-spectral-sequence-a.p001
Throughout this section, fix a commutative ring $\Bbbk$, a group $G$, and a
normal subgroup $H\lhd G$. Recall the restriction functor
$\Res^G_H:G\dcate{Mod}\to H\dcate{Mod}$ and the inflation functor
$\mathrm{Infl}^G_{G/H}:G/H\dcate{Mod}\to G\dcate{Mod}$ from
Definition~\ref{def:Res-Infl}. For every $G$-module $M$, introduce the
notation
\[ M^{H\lhd G}:=(\Res^G_HM)^H,\qquad
	M_{H\lhd G}:=(\Res^G_HM)_H. \]
Notice that $M^{H\lhd G}$ has an evident left $G/H$-action,
$(gH)\cdot x:=gx$. Similarly, $M_{H\lhd G}$ has an evident left
$G/H$-action. This gives two functors
\[ (\cdot)^{H\lhd G},\ (\cdot)_{H\lhd G}:
	G\dcate{Mod}\to G/H\dcate{Mod}. \]

% segment-id: o014.li2.chapter6.lhs-spectral-sequence-a.l001
\begin{lemma}
	The functor $(\cdot)^{H\lhd G}$ is right adjoint to
	$\mathrm{Infl}^G_{G/H}$, while $(\cdot)_{H\lhd G}$ is left adjoint to it.
\end{lemma}
\begin{proof}
	This is a straightforward extension of
	Proposition~\ref{prop:inv-coinv-infl}; verification is left to the reader.
\end{proof}

% segment-id: o014.li2.chapter6.lhs-spectral-sequence-a.p002
There are evident commutative diagrams of functors
\[\begin{tikzcd}[column sep=tiny]
	G\dcate{Mod} \arrow[rr,"{(\cdot)^{H\lhd G}}"]
	\arrow[rd,"{(\cdot)^G}"'] & &
	G/H\dcate{Mod} \arrow[ld,"{(\cdot)^{G/H}}"]\\
	& \Bbbk\dcate{Mod} &
\end{tikzcd}\quad\begin{tikzcd}[column sep=tiny]
	G\dcate{Mod} \arrow[rr,"{(\cdot)_{H\lhd G}}"]
	\arrow[rd,"{(\cdot)_G}"'] & &
	G/H\dcate{Mod}. \arrow[ld,"{(\cdot)_{G/H}}"]\\
	& \Bbbk\dcate{Mod} &
\end{tikzcd}\]

% segment-id: o014.li2.chapter6.lhs-spectral-sequence-a.p003
The functor $(\cdot)^{H\lhd G}$ is left exact and preserves injective
objects, since it has the exact left adjoint $\mathrm{Infl}^G_{G/H}$.
Similarly, $(\cdot)_{H\lhd G}$ is a right-exact functor that preserves
projective objects. Theorem~\ref{prop:derived-composite} on deriving composite
functors therefore gives canonical isomorphisms in the derived category
$\cate{D}(\Bbbk\dcate{Mod})$ (with appropriate parentheses omitted):
\begin{equation}\label{eqn:LHS-derived}
	\begin{gathered}
		\mathrm{R}(\cdot)^G\rightiso
		\mathrm{R}(\cdot)^{G/H}\circ\mathrm{R}(\cdot)^{H\lhd G},\\
		\mathrm{L}(\cdot)_G\leftiso
		\mathrm{L}(\cdot)_{G/H}\circ\mathrm{L}(\cdot)_{H\lhd G}.
	\end{gathered}
\end{equation}

% segment-id: o014.li2.chapter6.lhs-spectral-sequence-a.p004
The derived functors $\mathrm{R}(\cdot)^{H\lhd G}$ and
$\mathrm{L}(\cdot)_{H\lhd G}$ in these formulas are not new objects. For
each of the commutative diagrams
\[\begin{tikzcd}
	G\dcate{Mod} \arrow[r,"{(\cdot)^{H\lhd G}}"]
	\arrow[d,"{\Res^G_H}"'] & G/H\dcate{Mod}
	\arrow[d,"{\text{forgetful}=\Res^{G/H}_{\{1\}}}"]\\
	H\dcate{Mod} \arrow[r,"{(\cdot)^H}"'] & \Bbbk\dcate{Mod}
\end{tikzcd}\quad\begin{tikzcd}
	G\dcate{Mod} \arrow[r,"{(\cdot)_{H\lhd G}}"]
	\arrow[d,"{\Res^G_H}"'] & G/H\dcate{Mod}
	\arrow[d,"\text{forgetful}"]\\
	H\dcate{Mod} \arrow[r,"{(\cdot)_H}"'] & \Bbbk\dcate{Mod},
\end{tikzcd}\]
derive both composites according to Theorem~\ref{prop:derived-composite}.
Since $\Res^G_H$ is exact and preserves both injective and projective objects
(Proposition~\ref{prop:Ind-preservation}), while the forgetful functor is
exact, that theorem gives
\begin{gather*}
	\text{forgetful}\circ\mathrm{R}(\cdot)^{H\lhd G}
	\leftiso\mathrm{R}(\text{composite functor})
	\rightiso\mathrm{R}(\cdot)^H\circ\Res^G_H,\\
	\text{forgetful}\circ\mathrm{L}(\cdot)_{H\lhd G}
	\rightiso\mathrm{L}(\text{composite functor})
	\leftiso\mathrm{L}(\cdot)_H\circ\Res^G_H.
\end{gather*}
In other words, after applying the forgetful functor, the right-derived
functor $\mathrm{R}(\cdot)^{H\lhd G}$ is the same as
$\mathrm{R}(\cdot)^H$, and the left-derived functor
$\mathrm{L}(\cdot)_{H\lhd G}$ is the same as $\mathrm{L}(\cdot)_H$.
It is therefore natural to write their cohomology and homology as
\begin{align*}
	\mathrm{R}^n(\cdot)^{H\lhd G}
	&=\Hm^n(H,\cdot)\quad\text{with its $G/H$-action},\\
	\mathrm{L}_n(\cdot)_{H\lhd G}
	&=\Hm_n(H,\cdot)\quad\text{with its $G/H$-action}.
\end{align*}
Notice that we have begun to omit the restriction functor $\Res$ from the
notation.

% segment-id: o014.li2.chapter6.lhs-spectral-sequence-a.p005
These actions are easy to describe. For cohomology, let $M$ be a $G$-module
and $n\in\Z_{\geq0}$. There are at least three viewpoints on the $G/H$-action:
\begin{enumerate}
	\item Take an injective resolution $0\to M\to I^0\to\cdots$. It is also an
	injective resolution of $H$-modules. The group $G/H$ acts naturally on the
	complex $\cdots\to I^{n,H}\to I^{n+1,H}\to\cdots$, and hence on
	$\Hm^n(H,M)$. This description comes directly from the abstract
	construction above.
	\item The group $G/H$ acts on the functor
	$(\cdot)^H:G\dcate{Mod}\to\Bbbk\dcate{Mod}$. The universal property of
	universal cohomological $\delta$-functors then makes it act on
	$\Hm^n(H,M)$, compatibly with long exact sequences.
	\item In Example~\ref{eg:change-of-group-conj}, every $t\in G/H$ acts on
	$\Hm^n(H,M)$ through $\Hm^n(c_t)$.
\end{enumerate}
The reader is invited to verify that these three descriptions agree. The key
is to reduce their comparison to the case $n=0$.

% segment-id: o014.li2.chapter6.lhs-spectral-sequence-a.p006
The next important step is to extract more concrete information from
\eqref{eqn:LHS-derived}. This involves the cohomological and homological
bigraded spectral sequences of Definition~\ref{def:graded-spectral-sequence},
together with the following preparation.

% segment-id: o014.li2.chapter6.lhs-spectral-sequence-a.l002
\begin{lemma}\label{prop:LHS-action-aux}
	For every $G$-module $M$ and every $n$, the canonical homomorphism
	$\Hm^n(G,M)\to\Hm^n(H,M)$ (see \eqref{eqn:change-of-group-coh}) factors
	through $\Hm^n(H,M)^{G/H}$. The canonical homomorphism
	$\Hm_n(H,M)\to\Hm_n(G,M)$ factors through
	$\Hm_n(H,M)_{G/H}$.
\end{lemma}
\begin{proof}
	We discuss cohomology. Take an injective resolution $M\to I$, with
	$I=[I^0\to I^1\to\cdots]$, and write
	$I^H:=[I^{0,H}\to\cdots]$. Since $I$ is also an injective resolution of
	$M$ as an $H$-module, the inclusion $I^G\hookrightarrow I^H$ induces
	$\Hm^n(G,M)\to\Hm^n(H,M)$.

	As expected, Remark~\ref{rem:change-of-group-coh} ensures that this map is
	precisely the canonical homomorphism in question. Now compare it with the
	discussion of the $G/H$-action above.
\end{proof}

% segment-id: o014.li2.chapter6.lhs-spectral-sequence-a.t001
\begin{theorem}[R.\ Lyndon, G.\ Hochschild--J.-P.\ Serre]
\label{prop:LHS-SS}
	\index{spectral sequence!Lyndon--Hochschild--Serre}
	Let $M$ be a $G$-module and $H\lhd G$.
	\begin{enumerate}[(i)]
		\item There are first-quadrant cohomological and homological spectral
		sequences
		\begin{align*}
			E_2^{p,q}&\simeq\Hm^p\left(G/H,\Hm^q(H,M)\right)
			\Rightarrow\Hm^{p+q}(G,M),\\
			E^2_{p,q}&\simeq\Hm_p\left(G/H,\Hm_q(H,M)\right)
			\Rightarrow\Hm_{p+q}(G,M).
		\end{align*}
		\item Let $n\in\Z_{\geq1}$ and suppose that
		$\Hm^i(H,M)=0$ for all $1\leq i<n$ in the cohomological case, or
		$\Hm_i(H,M)=0$ for all $1\leq i<n$ in the homological case. Then there
		are, respectively, exact sequences
		\begin{equation*}\begin{split}
			0\to\Hm^n(G/H,M^H)&\to\Hm^n(G,M)\to\Hm^n(H,M)^{G/H}\\
			&\to\Hm^{n+1}(G/H,M^H)\to\Hm^{n+1}(G,M),
		\end{split}\end{equation*}
		and
		\begin{equation*}\begin{split}
			\Hm_{n+1}(G,M)&\to\Hm_{n+1}(G/H,M_H)\\
			&\to\Hm_n(H,M)_{G/H}\to\Hm_n(G,M)
			\to\Hm_n(G/H,M_H)\to0.
		\end{split}\end{equation*}
		Notice that for $n=1$ the hypothesis is vacuous.
		\item In the cohomological exact sequence above, every morphism other
		than $\Hm^n(H,M)^{G/H}\to\Hm^{n+1}(G/H,M^H)$ comes from the equivariant
		homomorphism $M^H\hookrightarrow M$ (relative to
		$G/H\twoheadleftarrow G$; see Definition~\ref{def:change-of-group-equiv}),
		the homomorphism $M\xrightarrow{\identity}M$ (relative to
		$G\hookleftarrow H$), and the factorization in
		Lemma~\ref{prop:LHS-action-aux}. The homological exact sequence is
		similar.
	\end{enumerate}
\end{theorem}
